\appendix
\section*{Appendix}

\section{Implementation Details}
\label{sec:implementation_details}


We implement \ours{} and all baselines on top of the \texttt{verl} fully-asynchronous pipeline, using a decoupled rollouter--trainer architecture with a staleness-bounded queue. We set \texttt{trigger\_parameter\_sync\_step} to $4$ and \texttt{staleness\_threshold} to $0.5$. We use AdamW with a learning rate $10^{-6}$, weight decay $0.1$, response length $4096$, rollout number $n=1$, global batch size 64, and epoch number 1. 

The maximum number of turns is $30$ for ALFWorld and ScienceWorld and $16$ for SearchQA.
%
The maximum generation length is $4096$ tokens per turn.
%
We use temperature $0.4$, top-$p=1.0$, and top-$k=-1$.
%
All evaluations use the ReAct prompt format with full chat history, no sliding window, thinking mode disabled, and \texttt{</action>} as the stop string.
%
For SearchQA, the agent is allowed at most $8$ search calls per episode and retrieves the top $3$ documents from a Wikipedia-2018 index using dense E5 retrieval with FAISS-Flat.
%
We evaluate on the standard ALFWorld valid-seen and valid-unseen splits, the ScienceWorld test set with $30$ task types and $5$ variations per task, and capped SearchQA splits with up to $500$ questions per split except Bamboogle, which contains $125$ questions.


For agentic OPD baselines, we keep the OPD backbone fixed and modify only the per-token or per-turn weighting scheme. TCOD-F2B truncates each rollout after
\begin{equation}
    k(n) = \min\left(k_{\mathrm{start}} + \frac{n}{\eta}, T\right)
\end{equation}
assistant turns. We set $k_{\mathrm{start}}=1$ and $\eta=1$. The maximum horizon is $T=30$ for ALFWorld and ScienceWorld, and $T=16$ for SearchQA, so the curriculum reaches the cap within one epoch. Entropy-Aware OPD scales the per-token weight by
\begin{equation}
    1 + \alpha_e \cdot \mathbf{1}
    \left[
    H_{\phi}(\cdot \mid h_t,a_{t,<j}) > \tau_e
    \right],
\end{equation}
where $H_{\phi}$ is estimated from the cached top-$K$ teacher distribution. We use $\tau_e=0.8$ and $\alpha_e=1.0$.

For \ours{}, the selective weight is defined as
\begin{equation}
    s_t = i_t \cdot c_t.
\end{equation}
The default per-turn intervention weight $i_t \in \{0,0.5,1\}$ ($\alpha=0.5$) is obtained from one short teacher query per assistant turn, using maximum generation length $1028$ and default temperature $0.0$. If the teacher output is unparseable, we conservatively fall back to $i_t=1$. We also retain the deterministic-failure short circuit: if the student turn is empty, unparseable, or violates the required action schema, we directly set $i_t=1$ without relying on the teacher intervention query.

We use the training split of each benchmark for training. For SearchQA \citep{jin2025search}, we randomly sample 10K examples from the training split as the training subset.

\paragraph{Evaluation Metrics.}
We use benchmark-specific metrics following the standard evaluation protocols.
%
For ALFWorld, we report success rate (SR) on the valid-seen and valid-unseen splits, where an episode is counted as successful if the environment returns \texttt{info["won"] = True} within the turn budget.
%
For ScienceWorld, we report both SR and mean score.
%
SR counts episodes whose cumulative score reaches $100$, while mean score averages the environment-provided cumulative reward over episodes.
%
Since ScienceWorld provides partial credit and small models rarely solve tasks end-to-end, mean score is a more sensitive measure of partial progress and catastrophic failures.
%
For SearchQA, we report exact match (EM), where the final answer is counted as correct if its normalized string matches any reference answer.
%
We follow the SearchR1 normalization protocol, which lowercases text, removes punctuation and articles, and collapses whitespace.
%
For SearchQA, we report EM on each evaluation split and macro-average across NQ, TriviaQA, PopQA, HotpotQA, 2WikiMultiHopQA, Musique, and Bamboogle.
%
We also report average number of turns for ALFWorld and ScienceWorld, and average search/turn statistics for SearchQA as auxiliary efficiency metrics.


\section{Supplementary Results}
\label{sec:supplementary_results}

\subsection{Case Study: Turn-Level Intervention Examples}
\label{sec:case_study}

To qualitatively examine how \ours{} allocates teacher supervision, we present representative ALFWorld cases for different intervention outcomes.
%
These examples illustrate that the intervention mechanism does not simply supervise every student response uniformly.
%
Instead, it distinguishes between acceptable actions, mildly suboptimal actions, clearly harmful actions, and deterministic execution failures.

\paragraph{Skip: acceptable student behavior.}
When the student response is reasonable under the current observation, \ours{} assigns \textsc{Skip} and suppresses the OPD loss for that turn.
%
For example, in a \texttt{pick\_two\_obj\_and\_place} task, the student observes the following environment state:
%
\begin{quote}
\small
\texttt{You arrive at coffeetable 1. On the coffeetable 1, you see a box 2, a creditcard 2, a creditcard 1, a houseplant 1, a plate 2, a plate 1, a statue 1, and a wateringcan 1.}
\end{quote}
%
The student responds:
%
\begin{quote}
\small
\texttt{Thought: I need to find two vases and place them on the coffeetable.}\\
\texttt{Action: look}
\end{quote}
%
The parsed action is \texttt{look}, and the teacher assigns \textsc{Skip} with intervention weight $i_t=0$.
%
Although the action does not immediately complete the task, it is a valid exploratory step under the current observation.
%
In another \texttt{pick\_clean\_then\_place\_in\_recep} task, the student similarly takes an exploratory action:
%
\begin{quote}
\small
\texttt{Thought: I need to examine the cabinet to find a knife.}\\
\texttt{Action: examine cabinet 1}
\end{quote}
%
The environment returns that the cabinet contains nothing, but the action itself is executable and plausible, so the teacher again assigns \textsc{Skip}.
%
These cases show that \ours{} avoids unnecessary token-level correction when the student is making plausible progress.

\paragraph{Weak: executable but suboptimal behavior.}
When the student action is valid but potentially inefficient or only partially aligned with the task, \ours{} assigns \textsc{Weak}.
%
For instance, in a \texttt{pick\_clean\_then\_place\_in\_recep} task, the student sees the same coffeetable observation and produces:
%
\begin{quote}
\small
\texttt{Thought: I need to find a clean kettle in the cabinet.}\\
\texttt{Action: examine cabinet 1}
\end{quote}
%
The action \texttt{examine cabinet 1} is executable and may still provide useful information, but it is not clearly the best next step for the task.
%
The teacher therefore assigns \textsc{Weak} with intervention weight $i_t=\alpha$.
%
As another example, in a \texttt{pick\_and\_place\_simple} task, the student reasons:
%
\begin{quote}
\small
\texttt{Thought: The cloth is not on the countertop. I'll check the cabinets to find a cloth.}\\
\texttt{Action: go to cabinet 1}
\end{quote}
%
The environment accepts the action and moves the agent to the closed cabinet, but the decision remains uncertain rather than clearly optimal.
%
The teacher again assigns \textsc{Weak}.
%
These examples demonstrate that \ours{} can provide partial supervision for ambiguous turns rather than forcing a binary choice between full imitation and no correction.

\paragraph{Strong: clearly harmful or inconsistent behavior.}
When the student response is clearly inconsistent with the task state or likely to push the trajectory in an unhelpful direction, \ours{} assigns \textsc{Strong}.
%
For example, in a \texttt{pick\_cool\_then\_place\_in\_recep} task, the student states that it should cool a plate in the fridge, but then navigates to a cabinet:
%
\begin{quote}
\small
\texttt{Thought: I need to find a plate and cool it in the fridge before putting it in the cabinet.}\\
\texttt{Action: go to cabinet 1}
\end{quote}
%
This response is internally inconsistent with the required cooling procedure, because moving to a cabinet does not help complete the cooling step.
%
The teacher assigns \textsc{Strong} with intervention weight $i_t=1$.
%
A related failure occurs when the student repeats the same high-level intent but selects another cabinet-navigation action:
%
\begin{quote}
\small
\texttt{Thought: I need to find a plate and cool it in the fridge before putting it in the cabinet.}\\
\texttt{Action: go to cabinet 10}
\end{quote}
%
Although the action is executable, it is unlikely to make progress toward cooling the object and may cause the trajectory to drift further from a successful plan.
%
These cases show that strong intervention is used when the current turn requires clear correction.

\paragraph{Deterministic failure: environment-detected execution error.}
In addition to teacher-judged intervention, \ours{} directly assigns full intervention when the environment or parser detects a deterministic failure.
%
For example, in a \texttt{pick\_and\_place\_simple} task, the student attempts:
%
\begin{quote}
\small
\texttt{Thought: I need to place the book on the sofa. I'll put the book on the sofa now.}\\
\texttt{Action: put book 1 in/on sofa 1}
\end{quote}
%
The parsed action is \texttt{put book 1 in/on sofa 1}, but the environment returns:
%
\begin{quote}
\small
\texttt{Nothing happens.}
\end{quote}
%
Similarly, in another episode, the student outputs:
%
\begin{quote}
\small
\texttt{Thought: I need to put the laptop in the safe. I'll put the laptop in the safe now.}\\
\texttt{Action: put laptop 1 in/on safe 1}
\end{quote}
%
Again, the environment returns \texttt{Nothing happens}, indicating that the action failed to produce a valid state transition.
%
In these cases, \ours{} bypasses semantic teacher judgment and directly applies full intervention with $i_t=1$.
%
This deterministic-failure pathway provides a high-precision correction signal for malformed, invalid, or ineffective executable actions.

\paragraph{Discussion.}
Overall, these case studies support the motivation for selective multi-turn OPD.
%
Different student turns require different levels of teacher supervision: some should be left unchanged, some benefit from weak guidance, and others require strong correction.
%
By using environment feedback together with teacher judgment, \ours{} allocates supervision at the turn level and avoids applying dense token-level OPD uniformly across the entire trajectory.
%
This behavior is especially important in multi-turn environments, where over-correcting acceptable exploration can reduce behavioral diversity, while failing to intervene on invalid or harmful actions can cause errors to compound in later turns.









\section{Prompt Templates}
\label{sec:appendix_prompts}


We use task-specific ReAct-style prompts for each benchmark. All prompts enforce a fixed output protocol so that the executable action can be reliably parsed and sent to the environment. We present the prompt details in \tabref{tab:prompt_alfworld}, \tabref{tab:prompt_scienceworld} and \tabref{tab:prompt_searchqa}. 


In addition to the task-specific agent prompts, \ours{} uses a separate intervention prompt to obtain turn-level teacher judgments.
%
Given the task goal, current environment state, admissible actions, and the student's candidate response, the teacher assigns one label from \{\textsc{Skip}, \textsc{Weak}, \textsc{Strong}\}.
%
The prompt is shown in \tabref{tab:prompt_intervention}.

\begin{table*}[t]
\centering
\caption{Prompt template for ALFWorld. We use a demonstration before the user turn.}
\label{tab:prompt_alfworld}
\resizebox{0.95\textwidth}{!}{%
\begin{tabular}{lp{0.78\textwidth}}
\toprule
\textbf{Field} & \textbf{Prompt Content} \\
\midrule
System prompt &
\begin{minipage}[t]{0.78\textwidth}
\ttfamily
You are an embodied agent solving a household task in a text-based simulator (AlfWorld). You receive an Observation describing what you can see and a list of Admissible commands you may execute. On every turn you must reply with exactly two lines: Thought: <one short sentence of reasoning> Action: <one command, copied verbatim from the Admissible list when possible> Plan briefly, then issue the next action. Only the line beginning with 'Action:' is used to step the environment. Do not add any other text after the Action line.
\end{minipage}
\\
\midrule
One-shot demonstration &
\begin{minipage}[t]{0.78\textwidth}
\ttfamily
An one-shot demonstration.
\end{minipage}
\\
\midrule
User turn &
\begin{minipage}[t]{0.78\textwidth}
\ttfamily
Task: \{task\_goal\}\\
Observation: \{scene\_description\}\\
Admissible: cmd1; cmd2; ...; cmd30 (+N more)
\end{minipage}
\\
\midrule
Assistant turn &
\begin{minipage}[t]{0.78\textwidth}
\ttfamily
Thought: <reasoning>\\
Action: <command>
\end{minipage}
\\
\bottomrule
\end{tabular}%
}
\end{table*}

\begin{table*}[t]
\centering
\caption{Prompt template for ScienceWorld. We use a single generic one-shot ReAct trajectory.}
\label{tab:prompt_scienceworld}
\resizebox{0.95\textwidth}{!}{%
\begin{tabular}{lp{0.78\textwidth}}
\toprule
\textbf{Field} & \textbf{Prompt Content} \\
\midrule
System prompt &
\begin{minipage}[t]{0.78\textwidth}
\ttfamily
You are an embodied agent solving a science task in the ScienceWorld text-based simulator. You receive an Observation describing what you can see and an Inventory listing what you carry. On every turn you must reply with exactly two lines: Thought: <one short sentence of reasoning> Action: <one valid action> Plan briefly, then issue the next action. Only the line beginning with 'Action:' is used to step the environment.
\end{minipage}
\\
\midrule
One-shot demonstration &
\begin{minipage}[t]{0.78\textwidth}
\ttfamily
A single generic ReAct trajectory.
\end{minipage}
\\
\midrule
User turn &
\begin{minipage}[t]{0.78\textwidth}
\ttfamily
Task: \{task\_description\}\\
Observation: \{observation\}\\
Inventory: \{inventory or "(empty)"\}
\end{minipage}
\\
\midrule
Assistant turn &
\begin{minipage}[t]{0.78\textwidth}
\ttfamily
Thought: <reasoning>\\
Action: <valid action>
\end{minipage}
\\
\bottomrule
\end{tabular}%
}
\end{table*}

\begin{table*}[t]
\centering
\caption{Prompt template for SearchQA. The agent repeatedly issues Wikipedia search queries and terminates with a concise final answer.}
\label{tab:prompt_searchqa}
\resizebox{0.95\textwidth}{!}{%
\begin{tabular}{lp{0.78\textwidth}}
\toprule
\textbf{Field} & \textbf{Prompt Content} \\
\midrule
System prompt &
\begin{minipage}[t]{0.78\textwidth}
\ttfamily
You are an expert open-domain question-answering agent. You answer questions by repeatedly issuing wikipedia search queries, reading the returned snippets, and finally producing a concise short answer.

Output protocol -- on every turn you must reply with exactly two blocks: <thought>... brief reasoning (1-2 sentences) ...</thought> <action>VERB[PAYLOAD]</action> where VERB[PAYLOAD] is exactly one of: search[your concise wikipedia query, 3-8 words] to fetch top-3 snippets answer[your final short answer] to terminate the episode Use search[...] until you have enough evidence, then answer[...]. Final answers should be short and direct (a single entity, date, or short phrase). Do not output anything after the </action> tag.
\end{minipage}
\\
\midrule
One-shot demonstration &
\begin{minipage}[t]{0.78\textwidth}
\ttfamily
A hardcoded one-shot demonstration.
\end{minipage}
\\
\midrule
User turn &
\begin{minipage}[t]{0.78\textwidth}
\ttfamily
Question: \{question\}? \\
\ttfamily
<information>1) ... 2) ... 3) ...</information>
\end{minipage}
\\
\midrule
Assistant turn &
\begin{minipage}[t]{0.78\textwidth}
\ttfamily
<thought>...</thought>\\
<action>search[query]</action>\\
\emph{or}\\
<thought>...</thought>\\
<action>answer[final]</action>
\end{minipage}
\\
\bottomrule
\end{tabular}%
}
\end{table*}


\begin{table*}[t]
\centering
\caption{Prompt template for turn-level intervention. The teacher acts as a supervision gatekeeper and assigns one intervention label for each student response.}
\label{tab:prompt_intervention}
\resizebox{0.95\textwidth}{!}{%
\begin{tabular}{lp{0.78\textwidth}}
\toprule
\textbf{Field} & \textbf{Prompt Content} \\
\midrule
System prompt &
\begin{minipage}[t]{0.78\textwidth}
\ttfamily
You are a SUPERVISION GATEKEEPER for distillation training of a smaller agent operating in a multi-turn text-based environment. Your job is to decide whether the student's proposed next action deserves teacher distillation supervision.
\end{minipage}
\\
\midrule
Episode goal &
\begin{minipage}[t]{0.78\textwidth}
\ttfamily
GOAL OF THE EPISODE:\\
Task type: \{task\_type\}\\
Task goal: \{task\_goal\}
\end{minipage}
\\
\midrule
Current environment state &
\begin{minipage}[t]{0.78\textwidth}
\ttfamily
CURRENT STATE OF THE ENVIRONMENT (what the student is reacting to):\\
Scene description: \{scene\_description\}\\
Last observation: \{last\_obs\}\\
Currently admissible actions (top \{admissible\_count\}):\\
\{admissible\}
\end{minipage}
\\
\midrule
Student response &
\begin{minipage}[t]{0.78\textwidth}
\ttfamily
STUDENT'S CANDIDATE NEXT ACTION (what we are evaluating):\\
\{llm\_text\}
\end{minipage}
\\
\midrule
Decision rubric &
\begin{minipage}[t]{0.78\textwidth}
\ttfamily
DECISION RUBRIC:\\
- SKIP: The student's action is reasonable, strategically equivalent to your preferred action, or the situation is too ambiguous for confident correction. Distillation here would just over-regularize the student.\\
- WEAK: The student's action may be suboptimal but the issue is uncertain or low-impact. A mild distillation signal is appropriate but the action is not clearly wrong.\\
- STRONG: The student's action is clearly problematic, off-task, redundant, or likely to fail the goal. The student needs to learn the correct behavior here.
\end{minipage}
\\
\midrule
Output format &
\begin{minipage}[t]{0.78\textwidth}
\ttfamily
Respond with EXACTLY ONE WORD: SKIP, WEAK, or STRONG.
\end{minipage}
\\
\bottomrule
\end{tabular}%
}
\end{table*}


\begin{algorithm}[t]
\caption{\ours{} Training Procedure}
\label{alg:sl_opd}
\begin{algorithmic}[1]
\REQUIRE Student $\pi_\theta$, teacher $\pi_\phi$, environment $\mathcal{E}$, horizon $T$, weak weight $\alpha$, constant $\epsilon$
\STATE Collect student on-policy trajectories $\tau \sim \pi_\theta$ by interacting with $\mathcal{E}$
\FOR{each turn $(h_t,a_t)$ in the collected trajectories}
    \IF{$a_t$ causes a deterministic execution failure}
        \STATE Set $z_t=\textsc{Strong}$ and $i_t=1$
    \ELSE
        \STATE Ask the teacher to assign $z_t \in \{\textsc{Skip},\textsc{Weak},\textsc{Strong}\}$
        \STATE Map the label to
        \[
        i_t =
        \begin{cases}
        0, & z_t=\textsc{Skip},\\
        \alpha, & z_t=\textsc{Weak},\\
        1, & z_t=\textsc{Strong}.
        \end{cases}
        \]
    \ENDIF
    \STATE Compute teacher confidence
    \[
    c_t =
    \frac{1}{|a_t|}
    \sum_{j=1}^{|a_t|}
    \max_v \pi_\phi(v \mid h_t,a_{t,<j}).
    \]
    \STATE Set the selective weight $s_t=i_t c_t$
\ENDFOR
\STATE Let $N=\sum_t |a_t|$ and compute
\[
Z=
\frac{N}{\max\left(\sum_t s_t |a_t|,\epsilon\right)}.
\]
\STATE Update the student using
\[
\mathcal{L}_{\mathrm{SL\mbox{-}OPD}}
=
Z\sum_t s_t \mathcal{L}_{\mathrm{OPD}}(h_t,a_t).
\]
\end{algorithmic}
\end{algorithm}